\chapter{Dark Web and the TOR Anonymous Network}

\setcounter{footnote}{73} % Continue numbering from Chapter 2

The Dark Web, through the anonymous network The Onion Router (TOR), is the environment in which this dissertation seeks to analyze the diffusion of power. The TOR network is one of the networks that compose the Dark Web, together with other networks that make an active effort to shield communications, such as the Freenet network and the I2P network. Specifically, the technical knowledge of the TOR network and, in general terms, the policies surrounding the anonymity conferred by this communication channel are essential for the analysis of power diffusion. Thus, this chapter grounds the reasons why we believe the Dark Web is the environment within the cyber domain that most challenges the authority of the State and emanates degrees of power to various actors -- whether public or private.

We have divided this second chapter into three main parts: the first part addresses the matters necessary to define and contextualize the Dark Web within the universe of the World Wide Web (WWW), part of cyberspace; the second part seeks, in the context of the Dark Web, to present the TOR anonymous network in its technical and political aspects; the third part establishes the TOR anonymous network as a relevant part of the knowledge structure in the 21st century.

In the first part, we will begin with the history of the global computer network, its origin and technical functioning. With this, we seek to show fundamental points for electronic devices to operate within the mesh of the Internet -- such as protocols and packet switching. The TCP/IP protocol, for example, was responsible for establishing connections between distinct computers, creating the system of digital information networks. Packet switching refers to the strategic addressing of data that travels across the network. However, it is quite common for individuals to use different technical terms to designate the same thing within the scope of the global computer network. For example, Finklea (2005) points out that users employ the term ``world wide web'' as a synonym for ``internet,'' although they are not the same. To prevent this from occurring, we felt the need to explain, beyond the history of the Internet, different technical terms such as ``network,'' ``Internet,'' ``WWW,'' ``Surface Web,'' ``Deep Web,'' ``Dark Web,'' and ``TOR'' -- in addition to protocols and packet switching. All of these terms are, in some sense, related to each other.

Next, we will present the WWW, commonly known as the ``web.'' The WWW is one of the application layers that operates through the Internet. And within this layer, according to current classification, there are six different categories of the Web: Surface Web, Opaque Web, Private Web, Proprietary Web, Truly Invisible Web, and Dark Web. The Surface Web is conventional navigation of the WWW (using browsers such as Google Chrome, Mozilla Firefox, Safari, Opera, etc.) through search portals like ``Google,'' ``Yahoo,'' and ``Bing!'' Unlike the Surface Web, the Deep Web is composed of the other five categories -- whose content is not indexed by the search engines of the aforementioned portals, meaning they will not appear in a list as results.

The second part aims to present the specific anonymous network for the purposes of analysis in this research -- The Onion Router (TOR). According to Finklea (2015), the TOR network, among the anonymous networks that compose the Dark Web, is the largest in terms of number of users and stored content. We will explain its origin, the ``onion routing'' technique -- responsible for encrypting messages and concealing the identities of users --, ``Tor Bridges,'' which assist in counter-censorship activities; the hidden services inherent to the network, its popularity, and its political aspect. Our objective was to make explicit the aspects of the anonymous network that make it a unique region of cyberspace.

In the third and final part, we sought to ground the reasons that lead us to believe that the TOR network is an important communication channel, and a relevant component of the knowledge structure in the 21st century.

Initially, the Internet was used by a technical elite composed of computer scientists and enthusiasts in the field. During this period, the moment prior to the creation of the ``WWW'' by Tim Berners-Lee, the navigation ``windows'' that exist today did not yet exist; to use the Internet, it was necessary to know a minimum of computer language. To create a more user-friendly interface for the Internet, and to solve problems of access to information, Tim Berners-Lee created the HTTP protocol, responsible for founding the World Wide Web (WWW), a system of connecting pages and servers using hyperlinks. In this environment, search engines emerged, such as Google, which navigate the WWW and store copies of pages and content, as well as addresses, on their own servers. On the search site, these engines are capable of showing lists of results based on terms entered in the search bar. Everything not catalogued by these major search engines remains ``hidden'' in the vast WWW, as they are not shown by these large information ``portals.'' This hidden content is called the Deep Web. It is in this region that the Dark Web is classified, and of which the TOR network is a part. Using specific software, it is possible to access this network. The non-state actors whose activities are studied in the final chapter of this dissertation use TOR to carry out their activities.

In other words, the understanding of the technical functioning of the Internet, in order to contextualize and explain the Dark Web and, specifically, the TOR anonymous network, is fundamental in this research for understanding how, and to what extent, the diffusion of power occurs in this region of cyberspace.

\section{History and Functioning of the Internet up to the World Wide Web}

When discussing the identity of the ``father of the Internet,'' many authors and journalists credit this title to the Englishman Tim Berners-Lee, creator of the World Wide Web -- a navigation system composed of pages and hyperlinks born in 1990. However, Berners-Lee's role was to create a navigation system with user-friendly interfaces, through texts and windows where one could ``jump'' between different pages with content. That is, the Internet already existed when he created the Web -- what the latter did was to allow greater reach of the global computer network by encompassing new users who lacked knowledge of computer language. (CASTELLS, 2001). For Naughton (1999), the title of ``father of the Internet'' is quite contested. But, according to him, Paul Baran has the strongest claims. A former employee of the Rand Corporation, Paul Baran was responsible for developing the basic principles behind the network system at the origin of the Internet: digital signals and data packets.\footnote{According to Naughton (1999), the Rand Corporation was a ``think tank'' founded primarily by the US Air Force which, in exchange for investments, received relative freedom to conduct research. Many of these major research projects were generally related to the Air Force's area of interest. When Paul Baran arrived at the Rand Corporation in 1959, at the height of the Cold War, he was tasked with investigating the possibility of preserving the systems responsible for strategic command and control weapons in case of physical attacks. There are other relevant names in the trajectory of Internet creation, besides Paul Baran. Banks (2008) recalls the role of Leonard Kleinrock, whose doctoral research on data distribution in a network resulted in a book published in 1964. Ceruzzi and Aspray (2008), in turn, attribute to Donald Davies, a British physicist, the development and coining of the term ``packet switching,'' essential for the construction of a distributed network.} At that time, the launch of the Soviet satellite Sputnik fueled the scientific and technological race between the USA and the USSR. (NAUGHTON, 1999). Faced with this scenario, the American Armed Forces invested in research that could develop technologies capable of protecting military communications against possible nuclear attacks.

In 1964, Paul Baran published his research on a communication system based on a distributive network capable of surviving physical attacks.\footnote{These first communication systems resistant to physical attacks were imagined by Baran in a way that they could resist without dependence on conventional electrical power: ``Since destruction of the national electricity grid was an assumed consequence of a nuclear exchange, the transmitter/receivers were to be powered by small generators fuelled by liquid petroleum from a 200-gallon tank buried in the ground beneath each tower. As each relay-tower would consume only fifty watts of electricity, this local supply of fuel would last for at least three months after failure of the grid.'' (NAUGHTON, 1999, p.105)} The survival of communications depended, broadly speaking, on the way in which the nodes (the term he used to designate communication stations) were connected: (a) centralized, with rare survival prospects; (b) decentralized, with low survival prospects; (c) arranged in a distributive network, with high survival prospects.\footnote{See ``Annex 3 -- Distributive Networks.''} (BARAN, 1964).

The next step would be to use digital signals to traverse the network, since analog signals lose their quality as they travel through the stations.\footnote{Digital signals do not degrade in the same way as analog signals because they are merely sequences of the digits ``zero'' and ``one'' -- in addition to there being sufficiently simple techniques to check whether a given sequence was transmitted correctly. (NAUGHTON, 1999).} According to Baran (1964), information would be transmitted between nodes using standardized message blocks. These blocks were of equal sizes and formed by data segments. To travel through the network to the final destination, the data blocks were ``disassembled'' so that their segments would traverse different routes. At the destination, these segments would be reordered to recreate the original message block, which would be read by the receiver. The nodes are of great importance because they would be responsible for choosing the best route for the next segment of the message block, so that all would arrive at the same chosen destination.\footnote{Naughton (1999, p.103) explains that, in order for the nodes to choose the best trajectory for each segment (lowest cost, highest speed), Baran created an algorithm containing a table. In this table was information about how many ``hops'' remained for the segment to reach the next node (neighboring node). The table thus indicated the best routes at any given time, as it was regularly fed with information about the state of neighboring nodes -- so that, if a node dropped out of this mesh, the table was updated to reflect the change and, thus, message blocks were rerouted to other paths.} The reason for all arriving at the same recipient is due to the message block itself: it contains information indicating the beginning of the message, addressing, sender, precedence, ``hand-over number,'' the actual message content in text format, and the end of the message. (BARAN, 1964).

An interesting metaphor that aids the understanding of this process is to imagine that the message, sent by the sender, is like a jigsaw puzzle. Sending the entire picture through the network (the information) proved to be a challenging task for the purposes of that era (maintaining communication channels intact in case of physical attack on infrastructure). It is easier to break up the picture so that, in the end, only the puzzle pieces (data segments) remain. Thus, they are sent individually through the network toward the designated receiver. The different pieces follow the most varied routes, at different speeds, constantly ``reading'' the table with information about the best route (lowest cost, highest speed). When they arrive at the destination, they are ``scrambled,'' but behind each piece there are instructions about its position to compose the image. The puzzle pieces are ``read'' and, in this way, it is possible to recreate the complete picture, just as at its origin. The final receiver can then observe the picture in its original format in the message read.\footnote{This system is somewhat different in the TOR anonymous network: message blocks and data segments also exist, however, the information behind each ``puzzle piece'' is hidden beneath layers of cryptography -- and each layer can only be ``read'' by a single node. There is no way to read all layers at once. Another difference concerns addressing. Each node would only know the address of the next node to which to send the piece, but not the final node. These and other characteristics compose the technology of the TOR network that differentiates it from the Internet network -- and, thus, attempts to protect messages and users. More on this in later sections.}

In sum, two points are central to Baran's work: that the necessary number of connections for each workstation (node) to ensure the survival of the entire network was three (this characteristic, the number of connections guaranteeing survival, he called ``redundancy''); and the segmentation of message blocks into smaller data, capable of being transmitted along the network to the final destination through digital signals and varied routes (this characteristic he termed ``packet switching'').

Baran's original work on packet switching of messages along a distributed network originated the fundamentals behind network communication that made possible the emergence of the Internet -- the network made of networks whose reach in the 1990s became global.
